\section{Emergent gauge dynamics: Maxwell and Yang--Mills}
\label{sec:gauge}

Section~\ref{sec:carrier} produced the algebraic curvatures $f_{\mu\nu}$ (trace) and
$\Gc_{\mu\nu}$ (traceless). We now show their dynamics are Maxwell and Yang--Mills.

\subsection{From the Brillouin zone to spacetime}
\label{subsec:bridge}

The Wilczek--Zee connection of Section~\ref{sec:carrier} lives over the Brillouin zone. The
physical gauge field lives over spacetime. The two are components of one connection on the
combined base the carrier depends on. Because the carrier is present throughout the emergent
spacetime, nothing forces its $U(3)$ frame to be the same everywhere; promoting the frame to a
slowly varying spacetime field $g(x)$ makes $B_\mu(x)\in\su{3}$ a genuine connection,
transforming as $B_\mu\to gB_\mu g^{-1}-\ii(\partial_\mu g)g^{-1}$. By the Wilczek--Zee
adiabatic theorem, the projected dynamics of the carrier amplitude $\psi\in\C^3$ is governed by
the covariant derivative $\Dc_\mu=\partial_\mu-\ii B_\mu$, so $B_\mu$ couples exactly as a gauge
potential, and the local $SU(3)$ acting on $(\psi,g)$ is an exact redundancy.

This is demonstrated, not merely argued \citep{repo}. Modelling slow spacetime dependence by a
non-rigid modulation of the carrier, the spacetime field strength $\Gc_{xx}(x)$ is nonzero and
$\su{3}$-valued ($\|\Gc_{\text{traceless}}\|\approx 0.08$), a genuine pure-gauge control
(constant carrier subspace, $U(3)$ basis rotation) gives $\|\Gc\|\approx 2.6\times10^{-9}$
confirming the field is physical, and the \emph{mixed} curvature $\Gc_{kx}$ is nonzero in every
$k$-direction. The Brillouin-zone object of Section~\ref{sec:carrier} and the spacetime gauge
field are therefore restrictions of a single $(k,x)$ connection, and the full $\su{3}$ structure
group certified over $k$ is the structure group of the spacetime field.

\subsection{Uniqueness of the effective actions}
\label{subsec:uniqueness}

The gauge field is a collective coordinate (the carrier frame); its entire dynamics is induced
by integrating out the gapped fluctuations. The leading effective action is fixed by symmetry.

\paragraph{Uniqueness statement.} Let $\Gamma[B]$ be the coarse-grained effective action for the
soft carrier-frame field, assuming (i) exact local $SU(3)$ (resp.\ $U(1)$) invariance, (ii)
locality (the integrated-out modes are gapped, so the kernel is short-ranged and admits a
derivative expansion), and (iii) Lorentz invariance in the emergent metric. Then its leading term
is uniquely
\begin{equation}
  \Gamma[B]=\int \dd^4x\,\sqrt{-g}\Bigl[-\tfrac{1}{4g^2}\,G^a_{\mu\nu}G^{a\,\mu\nu}\Bigr]+O(1/\Lambda^2),
  \qquad
  G^a_{\mu\nu}=\partial_\mu B^a_\nu-\partial_\nu B^a_\mu+f^{abc}B^b_\mu B^c_\nu,
  \label{eq:ym}
\end{equation}
and analogously $-\tfrac{1}{4e^2}f_{\mu\nu}f^{\mu\nu}$ for the abelian trace. By gauge invariance, $\Gamma$ depends only on $G_{\mu\nu}$ (resp.\ $f_{\mu\nu}$) and
covariant derivatives. Counting mass dimension ($[B]=1$, $[G]=2$): the dimension-2 mass term
$\Tr(B_\mu B^\mu)$ is \emph{not} gauge invariant and is forbidden (a massless gluon/photon);
the dimension-4 $\Tr(G_{\mu\nu}G^{\mu\nu})$ is the unique invariant (the topological
$G\tilde G$ is a total derivative); higher terms are irrelevant, suppressed by the gap scale
$\Lambda\sim 1/a$. Masslessness and the $G^2$ form are the same consequence of gauge invariance:
a pure-gauge configuration $B_\mu=-\ii(\partial_\mu g)g^{-1}$ has $G_{\mu\nu}=0$ (verified to
$2.3\times10^{-10}$) and hence zero action, which a mass term would penalize
\citep{repo}. The non-abelian structure constants are inherited from Section~\ref{sec:carrier}.

\subsection{Substrate-induced couplings}
\label{subsec:couplings}

The induced coupling is the substrate's own quantum-geometric response of the isolated carrier,
the integrated non-abelian quantum metric \citep{ProvostVallee1980,XiaoChangNiu2010}
\begin{equation}
  g_{\mu\nu}(\veck)=\tfrac12\operatorname{Re}\Tr\bigl[(\partial_\mu P_3)(\partial_\nu P_3)\bigr],
  \qquad \tfrac{1}{g^2}\propto\int_{\text{BZ}}\!\dd^4k\,\textstyle\sum_\mu g_{\mu\mu}(\veck)\ge 0,
\end{equation}
the symmetric partner of the Berry curvature. It is manifestly non-negative; numerically
$1/g^2$ is positive and finite ($\approx 12$, minimum sample $>0$), giving a \emph{stable
propagating} gauge field, and similarly $1/e^2>0$ for the abelian sector \citep{repo}. The
coefficient depends on strain/scale (it varies with the carrier amplitude), so the coupling
runs---weak and long-range in weak regions, strong and short-range where the link frames twist
hard, in the spirit of asymptotic freedom. The absolute normalization involves the
coarse-graining scheme; positivity, finiteness, and scale dependence are the theorem-level facts.

\subsection{Faraday $U(1)$ and charge}
\label{subsec:u1}

The abelian sector is the trace holonomy, $a_\mu=\tfrac13\Tr\Ac_\mu$, or the Berry connection of
any non-degenerate band. The Bianchi identity $\partial_{[\lambda}f_{\mu\nu]}=0$ is automatic
($f=\dd a$; verified to $10^{-13}$), giving the homogeneous Maxwell equations; the uniqueness
theorem gives the inhomogeneous $\partial^\mu f_{\mu\nu}=e^2 J_\nu$ with $\partial^\nu J_\nu=0$.
Electric charge is topological, $\pi_1(U(1))=\Z$: a carrier phase field with winding $n$ has
$\oint\partial_\mu\theta\,\dd x^\mu=2\pi n$ (verified for $n=0,\pm1,2,3$), so by Gauss's law the
enclosed charge is integer-quantized and conserved. The photon itself is a continuous,
non-quantized classical mode; any $h\nu$ granularity would enter only through coupling to
quantized matter (deferred).

A clean feature emerges: for the symmetric carrier the vacuum is \emph{EM-flat}---all abelian
Berry curvatures vanish pointwise ($\sim10^{-10}$ by the Wilson-loop measure), so empty space
carries no background electromagnetic field, as it should, while the abelian quantum metric
(hence $1/e^2$) is nonzero so the photon still propagates. Breaking the reality symmetry (a
generic deformation, i.e.\ matter) switches the abelian curvature on ($\sim 34$), so EM is
matter-sourced, consistent with charge-as-winding. This is fully compatible with the nonzero
$\su{3}$, which lives in the off-diagonal non-abelian sector \citep{repo}.
